\section{A dimension--spectral dispatch on a supplied face}
\label{sec:supplied-face-dispatch}

The preceding bridge identifies the linear system on a correct positive face.
There is then a useful choice that is independent of bipartiteness: exact CG
has a finite-dimensional cap, whereas interval Chebyshev iteration has the
usual condition-number rate.  Taking the better of these two bounds removes
an unnecessary logarithm whenever the supplied face is small.

Let $U\subseteq V$ be nonempty, put $n_U:=|U|$, and let
$D_U:=\diag(d_i:i\in U)$.  For the unregularized problem set $\rho=0$; for
the RPPR problem suppose that $U$ is the correct positive face and set
\begin{equation}
 c_U:=b_U-\alpha\rho D_U^{1/2}\mathbf 1_U.
 \label{eq:supplied-face-load}
\end{equation}
In either case the supplied-face solution is the unique vector
\begin{equation}
 Q_{UU}x_U^\dagger=c_U.
 \label{eq:supplied-face-system}
\end{equation}
The face is held fixed throughout this section.  A multiplication by
$Q_{UU}$ scans the adjacency list of every vertex in $U$, ignoring entries
outside $U$, and therefore costs $O(\operatorname{vol}(U))$.  Every dot
product, vector pass, and output write costs $O(n_U)$ and is absorbed because
$n_U\leq\operatorname{vol}(U)$.

For an iterate $x$ define the residual and its degree-normalized version by
\begin{equation}
 r(x):=c_U-Q_{UU}x,
 \qquad
 \widehat r(x):=D_U^{-1/2}r(x).
 \label{eq:supplied-face-residual}
\end{equation}
The next lemma records why a two-sided residual test is a safe semantic
certificate, even though a raw polynomial iterate need not be a nonnegative
lower envelope.

\begin{lemma}[Two-sided residual implies semantic accuracy]
\label{lem:supplied-face-two-sided-stop}
For every $x\in\mathbb R^U$,
\begin{equation}
 \|D_U^{-1/2}(x-x_U^\dagger)\|_\infty
 \leq \frac1\alpha\|\widehat r(x)\|_\infty.
 \label{eq:supplied-face-residual-semantic}
\end{equation}
Consequently, the computable two-sided stop
\begin{equation}
 \boxed{\ \|\widehat r(x)\|_\infty\leq\alpha\tau\ }
 \label{eq:supplied-face-two-sided-certificate}
\end{equation}
implies semantic error at most $\tau$.
\end{lemma}

\begin{proof}
Define the row-coordinate principal operator
\begin{equation*}
 H_U:=D_U^{-1/2}Q_{UU}D_U^{1/2}.
\end{equation*}
It is similar to the SPD Stieltjes matrix $Q_{UU}$, hence is a nonsingular
$M$-matrix and $H_U^{-1}\geq0$ entrywise.  On the full graph,
$QD^{1/2}\mathbf1=\alpha D^{1/2}\mathbf1$.  Passing to a principal face
removes only nonpositive off-diagonal terms, so
\begin{equation*}
 H_U\mathbf1_U\geq\alpha\mathbf1_U.
\end{equation*}
Multiplication by the nonnegative inverse gives
$H_U^{-1}\mathbf1_U\leq\alpha^{-1}\mathbf1_U$.  Therefore
$\|H_U^{-1}\|_{\infty\to\infty}\leq1/\alpha$.  Since
\begin{equation*}
 D_U^{-1/2}(x_U^\dagger-x)=H_U^{-1}\widehat r(x),
\end{equation*}
the claimed inequalities follow.
\end{proof}

For the sharp zero-start normalization used below, suppose that $s=e_v$,
$v\in U$, $0\leq\rho\leq1$, and $\rho\operatorname{vol}(U)\leq1$.  The last
condition is automatic on the correct RPPR support.  From
\eqref{eq:supplied-face-load},
\begin{align}
 \frac{\|c_U\|_2^2}{\alpha^2}
 &=\frac{(1-\rho d_v)^2}{d_v}
   +\rho^2\bigl(\operatorname{vol}(U)-d_v\bigr)\notag\\
 &\leq1+\rho\leq2.
 \label{eq:supplied-face-load-bound}
\end{align}
Here $\rho d_v\leq1$, so the first term is at most $1/d_v\leq1$, while the
second is at most $\rho^2\operatorname{vol}(U)\leq\rho$.  When $\rho=0$ the
same display gives the sharper $\|c_U\|_2\leq\alpha$ without a volume
assumption.

\begin{theorem}[Supplied-face CG/Chebyshev dispatch]
\label{thm:supplied-face-dispatch}
Assume $0<\alpha<1$, $0<\tau\leq1$, and the single-seed normalization and
load conditions of \eqref{eq:supplied-face-load-bound}.  Put
\begin{equation}
 \chi_\alpha:=\frac{1-\sqrt\alpha}{1+\sqrt\alpha},
 \qquad
 k_{\rm Ch}(\tau):=
 \left\lceil
  \frac{\log(2\sqrt2/\tau)}{-\log\chi_\alpha}
 \right\rceil.
 \label{eq:supplied-face-chebyshev-count}
\end{equation}
Starting from zero, one may obtain an iterate satisfying the stronger stop
\eqref{eq:supplied-face-two-sided-certificate} in at most
\begin{equation}
 N_U(\tau)\leq\min\{n_U,k_{\rm Ch}(\tau)\}
 \label{eq:supplied-face-dispatch-count}
\end{equation}
full principal-matrix products, by selecting exact-arithmetic CG when
$n_U\leq k_{\rm Ch}(\tau)$ and interval Chebyshev semi-iteration otherwise.
The fully charged real-arithmetic work, including a constant number of vector
passes per iteration and output, is
\begin{align}
 W_U
 &=O\!\left(\operatorname{vol}(U)
       \min\{n_U,k_{\rm Ch}(\tau)\}\right)\notag\\
 &=O\!\left(\operatorname{vol}(U)
       \min\left\{n_U,
       \alpha^{-1/2}\log\frac2\tau\right\}\right).
 \label{eq:supplied-face-dispatch-work}
\end{align}
For the accuracy regime $0<\tau\leq1/2$, the last logarithm is
$\Theta(\log(1/\tau))$.
For $\alpha=1$, $Q_{UU}=I$ and one face pass solves the system exactly.
\end{theorem}

\begin{proof}
The principal matrix has spectrum in $[\alpha,1]$.  Define the degree-$k$
residual polynomial
\begin{equation}
 p_k(\lambda):=
 \frac{T_k\!\left((1+\alpha-2\lambda)/(1-\alpha)\right)}
      {T_k\!\left((1+\alpha)/(1-\alpha)\right)},
 \label{eq:supplied-face-chebyshev-polynomial}
\end{equation}
where $T_k$ is the first-kind Chebyshev polynomial.  It satisfies $p_k(0)=1$
and the exact interval bound
\begin{equation}
 \max_{\lambda\in[\alpha,1]}|p_k(\lambda)|
 =\frac{2\chi_\alpha^k}{1+\chi_\alpha^{2k}}
 \leq2\chi_\alpha^k.
 \label{eq:supplied-face-chebyshev-factor}
\end{equation}
Chebyshev semi-iteration realizes the residual
$r_k=p_k(Q_{UU})c_U$ with $k$ principal products.  Since every degree is at
least one, \eqref{eq:supplied-face-load-bound} gives
\begin{equation}
 \|\widehat r_k\|_\infty
 \leq\|r_k\|_2
 \leq2\sqrt2\,\alpha\chi_\alpha^k.
 \label{eq:supplied-face-chebyshev-residual}
\end{equation}
Thus $k=k_{\rm Ch}(\tau)$ proves the two-sided stop.  Moreover,
\begin{equation}
 -\log\chi_\alpha
 =2\operatorname{artanh}(\sqrt\alpha)
 \geq2\sqrt\alpha,
 \label{eq:supplied-face-chebyshev-exponent}
\end{equation}
which proves the second line of \eqref{eq:supplied-face-dispatch-work}.

For the dimension cap, let $\lambda_1,\ldots,\lambda_s$ be the distinct
eigenvalues of $Q_{UU}$ that occur in the spectral support of $c_U$; then
$s\leq n_U$.  The error-annihilating polynomial
\begin{equation}
 a_s(z):=\prod_{h=1}^s\left(1-\frac z{\lambda_h}\right)
 \label{eq:supplied-face-annihilator}
\end{equation}
has degree $s$, satisfies $a_s(0)=1$, and annihilates $c_U$ under
$Q_{UU}$.  Hence
\begin{equation}
 q_{s-1}(z):=\frac{1-a_s(z)}z,
 \qquad
 x_U^\dagger=q_{s-1}(Q_{UU})c_U
 \in\mathcal K_s(Q_{UU},c_U).
 \label{eq:supplied-face-krylov-exactness}
\end{equation}
Exact-arithmetic CG therefore terminates with zero residual in at most
$s\leq n_U$ iterations.  Selecting the smaller certified count proves
\eqref{eq:supplied-face-dispatch-count}.  A sparse product and the constant
number of vector passes in either recurrence cost
$O(\operatorname{vol}(U))$; the final $O(n_U)$ output is absorbed.
\end{proof}

\begin{corollary}[Supplied correct-face PPR approximation]
\label{cor:supplied-face-ppr-dispatch}
Let $0<\alpha\leq1$ and $0<\varepsilon_{\rm ppr}\leq1$, set
$\rho=\tau=\varepsilon_{\rm ppr}/2$, and suppose that the exact RPPR support
$U=S^\star(\rho)$ is supplied.  Zero-pad the iterate returned by
\Cref{thm:supplied-face-dispatch} when $U\ne\varnothing$, and return zero
when $U=\varnothing$.  The resulting vector satisfies
\begin{equation}
 \|D^{-1/2}(x-x^0)\|_\infty\leq\varepsilon_{\rm ppr}
 \label{eq:supplied-face-ppr-error}
\end{equation}
in
\begin{equation}
 O\!\left(\operatorname{vol}(U)
 \min\left\{|U|,\alpha^{-1/2}
 \log\frac2{\varepsilon_{\rm ppr}}\right\}\right)
 \label{eq:supplied-face-ppr-work}
\end{equation}
real-arithmetic work, excluding the cost of supplying or independently
validating $U$.
\end{corollary}

\begin{proof}
If $U=\varnothing$, the returned zero vector equals $x^\star(\rho)$, so the
bias bridge alone gives error at most $\rho\leq\varepsilon_{\rm ppr}$.
Suppose now that $U\ne\varnothing$.  Then $v\in U$: otherwise every entry of
the shifted load \eqref{eq:supplied-face-load} is strictly negative, and the
entrywise nonnegative inverse of $Q_{UU}$ would give
$x_U^\star=Q_{UU}^{-1}c_U<0$, contradicting that $U$ is the positive support.
If $\alpha=1$, then $Q_{UU}=I$, so the one-pass identity solve is exact and
the bias bridge gives the claim.  Hence suppose $0<\alpha<1$.
The support-volume theorem gives
$\rho\operatorname{vol}(U)\leq1$, so
\Cref{thm:supplied-face-dispatch} applies.  Its two-sided stop and
\Cref{lem:supplied-face-two-sided-stop} give RPPR semantic error at most
$\tau$.  The bias bridge \Cref{thm:rppr-ppr-bias} contributes at most $\rho$.
Thus \eqref{eq:supplied-face-ppr-error} follows from $\rho+\tau$; substituting
$\tau=\varepsilon_{\rm ppr}/2$ gives \eqref{eq:supplied-face-ppr-work}.
For $\varepsilon_{\rm ppr}\leq1/2$, this is equivalently the usual
$O(\operatorname{vol}(U)\min\{|U|,
\alpha^{-1/2}\log(1/\varepsilon_{\rm ppr})\})$ notation.
\end{proof}

The constant $2\sqrt2$ comes only from the safe single-seed RPPR load.  For a
general nonzero supplied right-hand side, define
\begin{equation}
 k_{\rm Ch}(c_U,\tau):=
 \left\lceil
  \frac{\left[\log\!\left(2\|c_U\|_2/(\alpha\tau)\right)\right]_+}
       {-\log\chi_\alpha}
 \right\rceil,
 \qquad [z]_+:=\max\{0,z\}.
 \label{eq:supplied-face-general-load}
\end{equation}
The same dispatch holds with
$\min\{n_U,k_{\rm Ch}(c_U,\tau)\}$; only the normalized load removes this
additional scale from the logarithm.
When $c_U=0$, the zero start is already exact.

\paragraph{Relation to the fixed-face SOR factor.}
On a complete connected bipartite face, $\sigma_U=1$ in
\eqref{eq:bipartite-cross-block}.  Writing $q=\sqrt\alpha$ gives
\begin{equation}
 t_U=\frac{2q}{1+q^2},
 \qquad
 \zeta_U=\left(\frac{1-q}{1+q}\right)^2=\chi_\alpha^2.
 \label{eq:supplied-face-chebyshev-sor-factor}
\end{equation}
Thus one optimal red--black SOR sweep has the same asymptotic scalar factor as
two interval-Chebyshev factors, modulo the two primitives' scan conventions.
The new small-face escape is CG's finite dimension cap, not a better interval
condition-number exponent.

\paragraph{Scope boundary.}
This theorem is a supplied, fixed-face linear-solver result.  It assumes that
the correct positive RPPR face is already known (or that an unregularized
principal problem is intentionally supplied), and it returns the restricted
solution.  Polynomial and CG iterates may be signed and need not obey the
one-sided lower-envelope invariants used by local admission algorithms; this
is precisely why the stop in
\eqref{eq:supplied-face-two-sided-certificate} is two-sided.  For RPPR, the
outside inequalities \eqref{eq:rppr-inactive-face-test} must still be checked
unless correctness of the face is part of the input.

If admissions change $U$, then both $Q_{UU}$ and $c_U$ change.  Restarting
the dispatch on faces $U_1,U_2,\ldots$ costs only the unamortized sum
\begin{equation}
 \sum_j\operatorname{vol}(U_j)
 \min\{|U_j|,k_{\rm Ch}(\tau_j)\},
 \label{eq:supplied-face-restart-sum}
\end{equation}
which can repeatedly charge old rows.  Neither exact finite termination nor
the interval rate supplies the face, preserves a local admission certificate,
or controls this sum.  The arithmetic needed for CG dot products is charged
above, but distributed reduction latency is a separate resource;
finite-precision loss of exact termination, bit complexity, boundary
discovery, and the project's separate response and reporting resources remain outside
\eqref{eq:supplied-face-dispatch-work}.  The recommended use is therefore a
terminal dispatch on a certified face: compare $|U|$ with
$k_{\rm Ch}(\tau)$ and choose the dimension cap only when it is smaller.
